\section{Methodology}
\label{sec:methodology}

This section describes the dataset and label-audit procedure
(Section~\ref{sec:data}), the model architecture
(Section~\ref{sec:model}), the training configuration
(Section~\ref{sec:training}), and the deployment pipeline
(Section~\ref{sec:deployment}).

% -------------------------------------------------------------------------
\subsection{Dataset}
\label{sec:data}
% -------------------------------------------------------------------------

We use the publicly available \emph{Vehicle Detection} dataset hosted on
Roboflow Universe (workspace \texttt{personal-project-kej16}, project
\texttt{vehicle-detection-vznzd-dkl8g}, version~1, exported
2025-10-12) \cite{roboflow_dataset}.
The dataset is released under a Creative Commons CC~BY~4.0 licence,
permitting research use with attribution.

\textbf{Composition.}
The dataset contains 9\,725 images annotated with three vehicle classes:
\textit{bus}, \textit{car}, and \textit{van}.
Images are split into training (9\,218 images), validation (287 images),
and test (220 images) sets, corresponding to approximately
94.8\,\%, 2.9\,\%, and 2.3\,\% of the total, respectively.
Full composition statistics are given in Table~\ref{tab:dataset}.

\begin{table}[htbp]
\centering
\caption{Dataset composition per split (source: label audit stage of
\texttt{trace\_experiments.py}).}
\label{tab:dataset}
\begin{tabular}{lrrrrrr}
\toprule
Split & Images & Instances & Bus & Car & Van & Inst/img \\
\midrule
Train  & 9\,218 & 18\,494 & 4\,476 & 10\,739 & 3\,279 & 2.006 \\
Valid  &   287  &    427  &   113  &    279  &    35  & 1.488 \\
Test   &   220  &    413  &   109  &    236  &    68  & 1.877 \\
\midrule
Total  & 9\,725 & 19\,334 & 4\,698 & 11\,254 & 3\,382 & -- \\
\bottomrule
\end{tabular}
\end{table}

The \textit{car} class dominates all splits (approximately 58\,\% of
instances), while \textit{van} is the least frequent (approximately 17\,\%).
This imbalance is relevant because class frequency correlates with detection
difficulty and with the nature of confusion errors
(Section~\ref{sec:results}).

\textbf{Image preprocessing.}
All images are resized to $640\times640$ pixels at inference time using
letterboxing (padding with grey) to preserve the aspect ratio.
No offline resizing is performed; images are stored and read at their
original resolution.

% -------------------------------------------------------------------------
\subsection{Label-Format Audit and Cleaning}
\label{sec:audit}
% -------------------------------------------------------------------------

The Roboflow export uses the YOLO label format, in which each annotation
row in a \texttt{.txt} file contains either a bounding-box record
(five whitespace-separated values: class index, normalised centre-x,
centre-y, width, height) or a polygon record (more than five values:
class index followed by alternating x and y coordinates).

Upon inspecting the training labels, we found that three label files
contain both formats in the same file — specifically, bounding-box rows
mixed with polygon rows that encode bus outlines.
In Ultralytics~8.x (version 8.4.53 used here), the label parser
classifies a file as a polygon annotation if \emph{any} row contains more
than six tokens, and then reinterprets every row accordingly.
Bounding-box rows in such files are silently converted into degenerate
polygon approximations, which produces incorrect boxes for those
instances without raising any warning.

To correct this, we implemented a label-audit stage
(\texttt{stage\_audit} in \texttt{trace\_experiments.py}) that:
\begin{enumerate}
  \item Scans all label files in every split.
  \item Classifies each row by token count (5 tokens~=~bbox;
        $>6$ tokens~=~polygon).
  \item Flags files that contain both types as \emph{mixed}.
  \item Writes a cleaned copy of every label file by converting
        polygon rows to their axis-aligned bounding-box approximation
        using the formula
        $cx = (\min x + \max x)/2$, $cy = (\min y + \max y)/2$,
        $w = \max x - \min x$, $h = \max y - \min y$.
\end{enumerate}

Table~\ref{tab:audit} summarises the audit results.
No mixed files were found in the validation or test splits;
three mixed files appeared in the training split, all in the
\textit{bus} class.
The cleaned dataset is used for all reported evaluations.

\begin{table}[htbp]
\centering
\caption{Label-format audit results per split.}
\label{tab:audit}
\begin{tabular}{lrrrr}
\toprule
Split & BBox rows & Polygon rows & Mixed files & BBox in mixed \\
\midrule
Train & 18\,413 & 81 & 3 & 3 \\
Valid &    408  & 19 & 0 & 0 \\
Test  &    399  & 14 & 0 & 0 \\
\bottomrule
\end{tabular}
\end{table}

% -------------------------------------------------------------------------
\subsection{Model Architecture}
\label{sec:model}
% -------------------------------------------------------------------------

We use \textbf{YOLO26m}, the medium-scale variant of the YOLO26 detector
released by Ultralytics \cite{yolo262023}.
YOLO26 adopts a dual-head assignment strategy that eliminates the need for
non-maximum suppression (NMS) at inference time \cite{nmsfree2022}.
The one-to-one head produces the final prediction; the one-to-many head
provides auxiliary supervision during training only and is discarded at
inference.
Distribution Focal Loss is removed in favour of a standard regression loss,
and the model is optimised with the MuSGD optimiser.

The medium variant has the following characteristics, as measured by our
experiment harness:

\begin{itemize}
  \item \textbf{Parameters:} 74.70\,M
  \item \textbf{Computational cost:} 21.78\,GFLOPs (at $640\times640$)
  \item \textbf{Weight file size:} 42.0\,MB (\texttt{yolo26m.pt})
\end{itemize}

We initialise training from the publicly available COCO pre-trained
weights (\texttt{yolo26m.pt}) and fine-tune on the vehicle dataset.
The deployed checkpoint (\texttt{app/models/best.pt}) is the
best-validation-mAP$_{50}$ checkpoint from the initial training run
(epoch 23 of 50, as recorded in the checkpoint metadata).

For comparison, the experiment harness also supports three additional
baselines under identical settings:
\textbf{YOLOv8m} \cite{jocher2023ultralytics}, \textbf{YOLO11m} \cite{ultralytics_yolo11},
and \textbf{RT-DETR-l} \cite{rtdetr2022}.
Results for these baselines are reported as [DATA DIBUTUHKAN] pending
the multi-seed training runs.

% -------------------------------------------------------------------------
\subsection{Training Configuration}
\label{sec:training}
% -------------------------------------------------------------------------

All models are trained with a shared base configuration derived from the
deployed checkpoint's \texttt{train\_args} and encoded in
\texttt{trace\_experiments.py} (\texttt{BASE\_TRAIN}):

\begin{itemize}
  \item Image size: $640\times640$
  \item Epochs: 50 (early stopping patience: 100)
  \item Batch size: 16
  \item Optimiser: \texttt{auto} (Ultralytics selects AdamW for YOLO,
        SGD for RT-DETR)
  \item Initial learning rate ($\eta_0$): 0.01; final LR fraction: 0.01
  \item Momentum: 0.937; weight decay: $5\times10^{-4}$
  \item Warm-up epochs: 3
  \item Close-mosaic (epochs before end): 10
  \item Automatic mixed precision (AMP): enabled
  \item Deterministic mode: enabled (for reproducibility)
  \item Number of data-loader workers: 8
\end{itemize}

\textbf{Augmentation policy.}
All baselines are trained with a photometric-only augmentation policy
(\texttt{AUG\_PHOTOMETRIC}):
HSV hue $\pm0.015$, saturation $\pm0.3$, value $\pm0.3$;
all geometric augmentations (rotation, translation, scale, shear,
mosaic, MixUp, CutMix, copy-paste, random erasing) disabled.
This conservative policy avoids geometric deformations that might
corrupt the already-noisy bus annotations.

An augmentation ablation (Table~T9, [DATA DIBUTUHKAN]) compares three
policies on YOLO26m:
(a)~photometric only (baseline),
(b)~Ultralytics default (which includes mosaic and geometric
    augmentation), and
(c)~photometric~+ light geometric
    (fliplr~=~0.5, translate~=~0.1, scale~=~0.25).

All training runs use three random seeds (0, 1, 2) and report
mean\,$\pm$\,SD across seeds.
Runs whose output directory already contains \texttt{test\_metrics.json}
are skipped automatically, enabling incremental execution.

% -------------------------------------------------------------------------
\subsection{Deployment Pipeline}
\label{sec:deployment}
% -------------------------------------------------------------------------

The trained detector is packaged in a Streamlit web application named
\textbf{VehicleScope~AI} (version 1.0.0, \texttt{app/main.py}).
The application accepts uploaded images, runs YOLO26m inference at a
confidence threshold of 0.40 and an input size of 640, and displays
annotated bounding boxes colour-coded by class:
red for bus, teal for car, and yellow for van.

Vehicle counts are posted to a Google Sheets spreadsheet via an n8n
webhook (\texttt{app/utils/n8n\_client.py}).
The webhook payload includes a UTC timestamp, per-class counts, the total
vehicle count, the mean detection confidence, and an alert flag that is
set when the total exceeds a configurable threshold (default: 10 vehicles).
If the total exceeds the threshold, the n8n workflow additionally sends a
Telegram notification.

Figure~\ref{fig:pipeline} illustrates the end-to-end pipeline.
The deployment stack runs entirely on the user's local machine (Apple M4 Pro,
macOS~15.3.1, 24\,GB RAM) using the MPS backend for GPU-accelerated
inference when available.
CPU-only latency is measured as a fallback and reported in
Table~T10 (Section~\ref{sec:results}).
